% ---------------------------------------------------------------------
%  Chapitre 1 — Introduction (version aperçu, intégralement rédigée)
% ---------------------------------------------------------------------
\chapter{Introduction}
\label{chap:intro}

\section{But du mémoire}
\label{sec:but}

Lors de mon stage en entreprise, j'ai eu pour consigne de créer un agent
d'intelligence artificielle capable de faire du trading, c'est-à-dire
d'acheter et de revendre des actions en cherchant à tirer profit de leurs
variations. J'ai alors décidé de donner à cet agent des capacités
d'analyse statistique. Après avoir construit le nécessaire pour qu'il
puisse passer des ordres et accéder aux données du marché, j'ai estimé
qu'il devait être capable d'évaluer les stratégies de trading qu'il
aurait établies, une stratégie étant une règle fixe qui décide quand
acheter un titre et quand le revendre.

Je présente dans ce rapport l'outil d'évaluation de stratégies que j'ai
donné à mon agent.

L'objectif de cet outil n'est pas de découvrir une stratégie rentable.
Il est de construire un enchaînement de méthodes capable de déterminer,
pour une stratégie quelconque qu'on lui soumet, si l'avantage qu'elle
affiche est réel ou s'il relève du hasard d'échantillonnage ou du biais
de marché. Autrement dit, la contribution revendiquée est
méthodologique~: la chaîne de tests, sa justification et la vérification
de ses conditions d'application.

Les dix stratégies d'analyse technique utilisées dans ce rapport ne
constituent donc pas l'objet d'étude, mais un banc d'essai. Elles
fournissent des jeux de données réels permettant de montrer que l'outil
discrimine correctement, qu'il sait rejeter ce qui ne tient pas et
retenir ce qui résiste à l'analyse. Les résultats obtenus sont présentés
à titre d'illustration du bon fonctionnement du protocole, et non comme
des recommandations d'investissement. Pour s'assurer que l'outil ne
rejette pas tout par excès de sévérité, on lui soumettra aussi un témoin
dont l'avantage est réel par construction.

\section{Les trois menaces qui pèsent sur le jugement}
\label{sec:menaces}

La première difficulté tient à la dérive du marché. Le premier réflexe
pour juger une stratégie est d'examiner si elle est rentable~: un gain
positif la ferait « marcher ». Ce raisonnement est pourtant trompeur,
car encore faut-il que ce gain dépasse ce qu'une entrée au hasard aurait
rapporté. Ici les stratégies sont appliquées aux titres du S\&P~500,
l'indice des cinq cents grandes capitalisations américaines, dont la
tendance est haussière sur la période étudiée. Cette dérive haussière
pose problème, car une stratégie peut réaliser des gains sans avoir
capté le moindre signal. C'est pourquoi on ne mesurera pas les gains,
mais l'avantage qu'apporte chaque stratégie par rapport au hasard.

La deuxième difficulté provient du nombre de tests. Plus on teste de
stratégies, plus le risque d'obtenir au moins un faux positif augmente.
Ce problème sera traité au chapitre~\ref{chap:juger} à l'aide de la
correction de Bonferroni.

La troisième, sans doute la plus délicate, est la dépendance des
données. De nombreuses méthodes statistiques exigent l'indépendance des
observations~; or ici les trades sont dépendants sous plusieurs aspects,
comme la période, le titre ou encore le secteur. Cette difficulté sera
mesurée au chapitre~\ref{chap:donnees}, puis neutralisée au
chapitre~\ref{chap:juger}.

On est forcé de prendre en compte ces trois difficultés, sous peine d'un
verdict trop optimiste~: c'est leur traitement conjoint qui fait de ce
protocole un véritable outil de validation.

\section{Démarche du mémoire}
\label{sec:planification}

L'étude suit le fil d'une enquête unique, du diagnostic au verdict. Le
chapitre~\ref{chap:donnees} présente d'abord la matière première, la
base des quarante mille trades et la construction de l'avantage mesuré,
puis dresse l'état des lieux de ce qui menace les tests, en mesurant les
trois canaux par lesquels les observations se ressemblent. Le
chapitre~\ref{chap:juger} constitue le c\oe ur du travail~: on y juge
les stratégies, du test de référence jusqu'aux corrections de la
dépendance, en resserrant la prudence à chaque étape, avant de relire le
verdict à la lumière du comportement du marché lui-même. Le
chapitre~\ref{chap:contexte} change ensuite de question~: il ne s'agit
plus de valider, mais de comparer les stratégies entre elles et de
décrire les situations dans lesquelles elles se déclenchent. Le
chapitre~\ref{chap:protocole} récapitule enfin l'outil sous la forme
d'un protocole réutilisable, en assume les limites et trace les étapes
suivantes, avant la conclusion. Le fil conducteur va ainsi du plus
fragile, la donnée et ses dépendances, au plus solide, le verdict, puis
l'outil qui permettra d'en rendre d'autres.
